\chapter{Leptospirosis}
\index{Leptospirosis}

\section*{Definition and Overview}
Leptospirosis is a globally significant zoonotic bacterial infection caused by pathogenic spirochaetes of the genus \textit{Leptospira}. It is transmitted from animals to humans, representing a major public health concern particularly in tropical, subtropical, and rural agricultural regions. The clinical presentation of leptospirosis is notoriously variable, ranging from a mild, self-limiting febrile illness resembling influenza to a severe, life-threatening systemic disease known as Weil's disease, which is characterised by acute kidney injury, hepatic dysfunction, jaundice, and severe haemorrhage. Early recognition and prompt antibiotic therapy are crucial to prevent progression to severe complications and death.

\section*{Epidemiology}
Leptospirosis occurs worldwide but is most prevalent in warm, humid climates where surface water and soil remain moist, facilitating bacterial survival. Globally, it is estimated to cause over one million clinical cases and approximately 60,000 deaths annually. In many countries, the disease is relatively uncommon but remains endemic, with the majority of cases reported in tropical north Queensland. It is primarily an occupational hazard, frequently affecting dairy farmers, banana growers, sugar cane cutters, abattoir workers, veterinarians, and military personnel. Outbreaks are strongly associated with heavy rainfall, flooding, and natural disasters, which wash animal urine into surrounding water systems and soil, exposing the general public.

\section*{Aetiology and Pathophysiology}
Leptospirosis is caused by spiral-shaped, highly motile bacteria of the genus \textit{Leptospira}. The pathogen's transmission and subsequent pathobiology include:
\begin{itemize}
    \item \textbf{Reservoir hosts:} pathogenic \textit{Leptospira} species chronically colonise the renal tubules of various wild and domestic animals, particularly rodents (the primary reservoir), cattle, pigs, dogs, and horses, which excrete the bacteria in their urine without showing signs of illness.
    \item \textbf{Transmission pathways:} humans contract the infection through direct contact with infected animal urine, or indirectly through contact with fresh water, wet soil, or vegetation contaminated by contaminated urine. The spirochaetes enter the body through cuts, abrasions, intact mucous membranes (such as the mouth, nose, and conjunctivae), or by inhaling aerosolised droplets.
    \item \textbf{Leptospiraemic phase:} once inside the host, the bacteria rapidly enter the bloodstream and spread throughout the body, localising in organs including the liver, kidneys, spleen, central nervous system, and skeletal muscles.
    \item \textbf{Endothelial damage and vasculitis:} the primary pathophysiological mechanism is a widespread immune-mediated systemic vasculitis. The bacteria damage the endothelial lining of small blood vessels, leading to capillary leakage, tissue oedema, and localised haemorrhage.
    \item \textbf{Organ-specific dysfunction:} renal involvement results in acute interstitial nephritis and tubular necrosis, causing acute kidney injury. Hepatic involvement causes hepatocellular dysfunction and jaundice without significant liver necrosis. Severe pulmonary haemorrhage syndrome represents an immune-mediated alveolar injury leading to massive bleeding in the lungs.
\end{itemize}

\section*{Clinical Features}
The incubation period of leptospirosis ranges from 2 to 30 days (typically 7 to 14 days). The clinical features are classically biphasic, starting with an acute septicaemic phase followed by an immune phase, and are categorised as follows:
\begin{itemize}
    \item \textbf{Anicteric leptospirosis (mild to moderate):} accounts for approximately 90\% of cases and presents with a sudden onset of high fever, chills, severe frontal headache, intense muscle aches (predominantly in the calves, thighs, and lower back), conjunctival suffusion (redness of the eyes without discharge), abdominal pain, nausea, vomiting, and diarrhoea. A transient macular skin rash may also appear.
    \item \textbf{Icteric leptospirosis (Weil's disease):} a severe presentation (occurring in about 10\% of cases) marked by deep jaundice, significant hepatic tenderness, oliguric or anuric renal failure, myocarditis, and a severe bleeding diathesis manifesting as epistaxis, haematemesis, melaena, or purpura.
    \item \textbf{Pulmonary haemorrhage syndrome:} characterised by acute respiratory distress, shortness of breath, a rapid respiratory rate, and coughing up blood (haemoptysis), which can lead to rapid asphyxiation and respiratory failure.
    \item \textbf{Aseptic meningitis:} a common feature of the immune phase, manifesting as a stiff neck, photophobia, severe headache, and vomiting.
\end{itemize}

\section*{Differential Diagnosis}
Because the early symptoms of leptospirosis are non-specific, it is frequently misdiagnosed. Clinicians must differentiate it from other acute febrile illnesses, including:
\begin{itemize}
    \item \textbf{Influenza:} distinguished by the presence of prominent upper respiratory tract symptoms, such as sore throat and rhinorrhoea, which are typically absent in leptospirosis.
    \item \textbf{Dengue fever:} characterised by severe retro-orbital pain, severe joint pain ("breakbone fever"), leucopenia, and a distinct maculopapular rash, though it shares thrombocytopenia with leptospirosis.
    \item \textbf{Malaria:} must be ruled out in travellers returning from endemic areas; diagnosed by detecting Plasmodium parasites on thick and thin blood films.
    \item \textbf{Rickettsial infections (such as scrub typhus):} often present with a characteristic black eschar at the site of the vector bite and local lymphadenopathy.
    \item \textbf{Viral hepatitis:} presents with jaundice and elevated liver transaminases, but lacks the severe muscle pain, conjunctival suffusion, and renal involvement characteristic of leptospirosis.
    \item \textbf{Bacterial meningitis:} presents with meningism, but requires cerebrospinal fluid analysis to differentiate from the aseptic meningitis of leptospirosis.
\end{itemize}

\section*{When to Seek Medical Care}
Anyone who develops a sudden febrile illness accompanied by severe headache and muscle aches, particularly after contact with rodents, farm animals, or soil/water contaminated by animal urine (especially after flooding or recreational freshwater activities), should seek prompt medical advice.

\textbf{Contact your local emergency services for:}
\begin{itemize}
    \item Difficulty breathing, a rapid breathing rate, or coughing up blood (haemoptysis)
    \item Jaundice (yellowing of the skin or the whites of the eyes)
    \item Severe reduction in urine output or complete inability to pass urine
    \item Confusion, severe drowsiness, seizures, or altered mental state
    \item Unexplained bleeding, such as frequent nosebleeds, vomiting blood, or severe bruising
\end{itemize}

\section*{Diagnosis and Investigations}
Confirming a diagnosis of leptospirosis requires a high index of clinical suspicion combined with specific microbiological and serological tests:
\begin{itemize}
    \item \textbf{Microscopic Agglutination Test (MAT):} the reference standard serological test, which detects agglutinating antibodies using live antigen suspensions; however, it requires acute and convalescent sera collected 10 to 14 days apart to demonstrate a fourfold rise in antibody titre.
    \item \textbf{Enzyme-Linked Immunosorbent Assay (ELISA):} detects Leptospira-specific IgM antibodies, which become detectable from the end of the first week of illness.
    \item \textbf{Polymerase Chain Reaction (PCR):} a highly sensitive molecular test that can detect Leptospira DNA in whole blood or cerebrospinal fluid (CSF) during the first 7 days of illness (septicaemic phase), and in urine after the first week.
    \item \textbf{Bacterial culture:} culturing the organism from blood or CSF during the first week of illness, or from urine thereafter; however, this is clinically limited due to the slow growth rate of the bacteria, which can take weeks.
    \item \textbf{Routine biochemistry and haematology:} full blood count typically reveals neutrophilic leucocytosis and thrombocytopenia. Kidney function tests show elevated creatinine and urea. Liver function tests reveal elevated bilirubin with mild to moderate elevations in transaminases. Urinalysis may show proteinuria, haematuria, and casts.
\end{itemize}

\section*{Management}
Management is directed at eliminating the bacterial pathogen and providing supportive therapy for organ dysfunction:
\begin{itemize}
    \item \textbf{Antibiotic therapy:} must be initiated as early as possible. Mild cases are treated with oral doxycycline (100 mg twice daily) or oral amoxicillin. Severe cases require intravenous benzylpenicillin, ceftriaxone, or cefotaxime.
    \item \textbf{Jarisch-Herxheimer reaction:} clinicians must monitor the patient closely for this reaction within the first 24 hours of starting antibiotics, characterised by a sudden onset of fever, chills, hypotension, and worsening symptoms due to the rapid release of bacterial endotoxins.
    \item \textbf{Renal support:} close monitoring of fluid balance is essential. Patients with acute kidney injury may require temporary renal replacement therapy (haemodialysis or haemofiltration).
    \item \textbf{Pulmonary and circulatory support:} severe cases with pulmonary haemorrhage require admission to an intensive care unit (ICU) for invasive mechanical ventilation and careful fluid management. Transfusions of platelets and packed red blood cells may be necessary to correct thrombocytopenia and blood loss.
\end{itemize}

\section*{Complications}
The complications of untreated or severe leptospirosis can be devastating and include:
\begin{itemize}
    \item Acute kidney injury requiring dialysis
    \item Severe pulmonary haemorrhage syndrome, which can lead to rapid respiratory failure and asphyxiation
    \item Liver failure with severe hyperbilirubinaemia (jaundice)
    \item Myocarditis, cardiac arrhythmias, and cardiovascular collapse
    \item Aseptic meningitis, encephalitis, and peripheral neuropathies
    \item Uveitis, a late inflammatory eye condition that can occur months after recovery
    \item Miscarriage, stillbirth, or congenital infection in pregnant individuals
\end{itemize}

\section*{Prognosis and Outcomes}
The prognosis for the majority of patients with mild, anicteric leptospirosis is excellent, with complete recovery and no long-term sequelae. However, the prognosis worsens significantly for patients who develop severe disease. The mortality rate for icteric leptospirosis (Weil's disease) ranges from 5\% to 15\%, and it can exceed 50\% in cases of severe pulmonary haemorrhage syndrome, even with access to modern intensive care. Recovery from severe disease can be prolonged, with some patients experiencing persistent fatigue, muscle weakness, headaches, and psychological symptoms for several months.

\section*{Prevention and Self-Care}
Preventative strategies focus on reducing exposure to infected animals and contaminated water or soil:
\begin{itemize}
    \item Avoid swimming, wading, or using recreational watercraft in lakes, rivers, or canals that may be contaminated with animal urine, particularly after periods of heavy rain or flooding
    \item Wear appropriate protective clothing, such as waterproof boots, gloves, long-sleeved shirts, and safety glasses, when working with farm animals, handling soil, or cleaning up flood-affected areas
    \item Control rodent populations in and around residential, agricultural, and workplace environments by using traps, baits, and clearing debris
    \item Vaccinate domestic livestock (cattle, pigs) and companion pets (dogs) against leptospirosis to reduce shedding of the bacteria
    \item Practice thorough hand hygiene by washing hands with soap and water after animal contact, gardening, or outdoor activities
    \item Cover any cuts, scrapes, or open wounds with clean, waterproof dressings before engaging in outdoor activities or working with animals
\end{itemize}

\section*{Sources and Further Reading}
The following reputable organisations provide detailed information on leptospirosis:
\begin{itemize}
    \item Healthdirect Australia. \textit{Leptospirosis overview, risks, and prevention.} https://www.healthdirect.gov.au/
    \item Queensland Health. \textit{Leptospirosis factsheet and clinical guidelines.} https://www.health.qld.gov.au/
    \item World Health Organization. \textit{Leptospirosis global burden and public health response.} https://www.who.int/
    \item Centers for Disease Control and Prevention. \textit{Leptospirosis transmission, symptoms, and diagnosis.} https://www.cdc.gov/
    \item Better Health Channel. \textit{Leptospirosis prevention and health information.} https://www.betterhealth.vic.gov.au/
\end{itemize}
